\chapter{Conclusiones y líneas futuras}\label{cap:conclusiones}

Este capítulo cierra el trabajo revisando en qué medida se han alcanzado los objetivos planteados en el Capítulo~\ref{cap:intro}, señalando las aportaciones que lo distinguen, reconociendo sus limitaciones y proyectando su continuidad. Conviene recordar el marco desde el que se leen estas conclusiones: se trata de un \textbf{Trabajo de Fin de Máster de diseño}, cuyo producto es una secuencia didáctica fundamentada y sus materiales, no un estudio empírico de aula. Por ello, los resultados que aquí se recogen se refieren al proceso de diseño, a la coherencia interna de la propuesta y a su alineación con la literatura y con el currículo; su eficacia sobre el aprendizaje es una hipótesis razonada que deberá contrastarse en la implementación prevista.

\section{Grado de consecución de los objetivos}

El \textbf{objetivo general} ---diseñar una secuencia didáctica basada en el modelo 5E para la determinación de la edad y el tamaño del universo, empleando la expansión y la ley de Hubble-Lemaître como vía de medida--- se ha cumplido: el trabajo entrega una secuencia completa para 2.º de Bachillerato que culmina, precisamente, en la estimación de esos dos valores a partir de observaciones reales. Los objetivos específicos se han abordado del modo siguiente:

\begin{itemize}
\item \textbf{OE1 y OE2} (Cap.~\ref{cap:revision}). La revisión sistematizada de la literatura sobre enseñanza de la cosmología moderna en secundaria se ha materializado en el análisis de 78 fuentes, organizadas en diez categorías de estrategia y contexto, con identificación de los contenidos más tratados, las estrategias empleadas y las dificultades conceptuales documentadas. De esa síntesis se extrajeron los elementos transferibles y se asignaron a las cinco fases del modelo 5E.

\item \textbf{OE3} (Cap.~\ref{cap:revision}). El cotejo de esos contenidos con el currículo de Física que imparte el autor situó la propuesta en Pearson Edexcel IAL (Topic~11, «Astrophysics and Cosmology»), que encaja casi punto por punto con el tema y fija el nivel en 2.º de Bachillerato; la PCE UNED se valoró y quedó descartada como marco de la propuesta.

\item \textbf{OE4} (Cap.~\ref{cap:propuesta}). Sobre esa base se diseñó la secuencia 5E, seleccionando qué contenidos incluir, qué estrategias resultaban más adecuadas para cada fase y qué ideas erróneas debían abordarse de forma explícita, con una estructura articulada en torno a una pregunta motriz y a una progresión de conceptos que conduce del corrimiento al rojo hasta la edad y el tamaño del universo.

\item \textbf{OE5} (Cap.~\ref{cap:desarrollo}). Se diseñaron los instrumentos de evaluación: dos rúbricas de cuatro niveles, un cuestionario conceptual pre-post adaptado de inventarios validados \parencite{gasparini2025} y un instrumento formativo integrado en la web interactiva (confrontación de ideas erróneas mediante verdadero/falso), previendo su aplicación en la implementación del curso 2026-2027.

\item \textbf{OE6} (Cap.~\ref{cap:desarrollo}). Se propuso una metodología didáctica para la lectura e interpretación de las ecuaciones que sustentan la medida ---no para su deducción formal---, articulada en seis etapas y aplicada al catálogo de ecuaciones de la secuencia, como aportación transversal recogida en el apartado correspondiente y en el anexo de ecuaciones.
\end{itemize}

En conjunto, los seis objetivos específicos se consideran alcanzados en el plano del diseño, que es el alcance declarado del trabajo.

\section{Aportaciones del trabajo}

Más allá de entregar una secuencia 5E coherente y fundamentada, el trabajo aporta dos elementos que lo singularizan.

El primero es la \textbf{web interactiva de acompañamiento \textit{Cosmos 5E}}: una herramienta offline, de código abierto y sin dependencias, que recorre las cinco fases y ofrece como piezas centrales una calculadora didáctica y un diagrama de Hubble construido sobre datos reales \parencite{steer2017}. Su principio de diseño es no funcionar como una caja negra: cada paso explica qué se va a hacer, por qué interviene y qué significa el resultado, mostrando la cadena de la medida con sus unidades desplegadas. El uso de distancias y velocidades reales, en lugar de valores fabricados, ancla la propuesta en la práctica científica auténtica y permite que el alumnado obtenga su propia estimación de la constante de Hubble y, a partir de ella, de la edad y el tamaño del universo.

El segundo es la \textbf{metodología para la lectura e interpretación de ecuaciones}, organizada en seis etapas (símbolos, relación, sentido, propósito, origen y tipo). Concebida como articulación operativa de principios ya establecidos ---la conversión entre registros de representación \parencite{duval2006}, la alfabetización matemática y la naturaleza de la ciencia---, ofrece un procedimiento fijo y repetible que convierte la comprensión de una fórmula en una tarea observable, que puede plantearse, guiarse y calificarse. Aunque aquí se aplica a la cosmología ---un banco de pruebas idóneo, por cuanto sus ecuaciones son imprescindibles pero indeducibles a este nivel---, la metodología es transferible a otros contenidos de la física.

\section{Limitaciones}

El trabajo presenta limitaciones que conviene explicitar y que, en su mayoría, derivan de su naturaleza de propuesta de diseño.

\begin{itemize}
\item \textbf{Ausencia de pilotaje.} Al no implementarse dentro del propio TFM, la eficacia de la secuencia sobre el aprendizaje no se ha comprobado empíricamente: es una hipótesis fundamentada, no un resultado.

\item \textbf{Contexto único.} El diseño se ancló en un currículo concreto (Edexcel IAL), lo que favorece su encaje pero limita la generalización directa a otros marcos, como la PCE UNED, que quedó fuera de la propuesta.

\item \textbf{Validación pendiente de la web y de los instrumentos.} La herramienta no se ha sometido a una evaluación de usabilidad y accesibilidad con alumnado real, y los instrumentos de evaluación, aunque adaptados de inventarios validados, no se han validado psicométricamente en el contexto de aplicación.

\item \textbf{Carácter didáctico del conjunto de datos.} El conjunto de galaxias del diagrama de Hubble es una selección con fines didácticos, no un muestreo científico; se ha priorizado la claridad del ajuste lineal y la aparición de las velocidades peculiares como lección, sobre la exhaustividad.

\item \textbf{Régimen simplificado.} La propuesta se mantiene deliberadamente en el régimen lineal de la ley de Hubble y en la aproximación de la edad como inversa de la constante de Hubble; los tratamientos exactos quedan señalados como matiz o como ampliación, no desarrollados.
\end{itemize}

\section{Líneas futuras}

La continuación natural del trabajo es su \textbf{implementación en el aula durante el curso 2026-2027}, en la ventana ya identificada del currículo de Física del autor, con la aplicación de los instrumentos diseñados en un esquema de evaluación \textbf{pre-post}. Esa experiencia permitiría contrastar la hipótesis central del trabajo y recoger evidencia sobre la comprensión efectiva de la edad y el tamaño del universo.

De ella se derivan otras líneas:

\begin{itemize}
\item \textbf{Validación y mejora de la web}, incorporando la observación del uso real del alumnado, criterios de accesibilidad y, en su caso, nuevas actividades o datos.

\item \textbf{Transferencia a otros contextos}, adaptando la secuencia a marcos curriculares distintos y extendiendo la metodología de lectura de ecuaciones a otros bloques de la física.

\item \textbf{Ampliación de nivel}, para itinerarios más avanzados, abordando las ecuaciones de Friedmann y la geometría del universo con mayor profundidad de la que admite el nivel de referencia.

\item \textbf{Integrar más actividades y elementos de evaluación en la web}: llevar a la herramienta actividades que hoy se apoyan en materiales complementarios y los elementos de evaluación ---registro de las respuestas, aplicación de las rúbricas y seguimiento del progreso del alumnado---, de modo que la secuencia completa y su evaluación puedan realizarse y quedar registradas en línea.

\item \textbf{Entorno docente en la web}: incorporar a la propia herramienta una sección docente ---con navegación libre por las fases, la guía de sesiones, los instrumentos y sus claves de corrección, y los datos descargables--- e integrar en la misma aplicación los dos modos de uso, el del alumnado y el del profesorado. De este modo, los materiales de reproducibilidad que este trabajo recoge en los anexos quedarían también accesibles en línea, sin necesidad de preparar nada por separado.
\end{itemize}

\section{Reflexión final}

La secuencia diseñada busca que el alumnado no reciba la edad y el tamaño del universo como cifras cerradas, sino que las reconstruya siguiendo el mismo camino que la ciencia: observar, medir, interpretar y, sobre todo, comprender qué significan los números y de dónde proceden. En coherencia con el espíritu del modelo 5E, el final de la secuencia se plantea como un principio: al término del recorrido, el estudiante debería quedarse no solo con respuestas, sino con mejores preguntas. Si la propuesta consigue que alguien mire el cielo y entienda que su oscuridad, su antigüedad y su extensión son cosas que se pueden medir y razonar, habrá cumplido su verdadero propósito.